% Generated draft: requires mathematical and source review.
\documentclass[11pt]{article}
\usepackage[margin=25mm]{geometry}
\usepackage{fontspec}
\setmainfont{Latin Modern Roman}
\usepackage{amsmath,amssymb,mathtools,unicode-math}
\setmathfont{Latin Modern Math}
\usepackage{enumitem,needspace,xurl,hyperref}
\setlength{\parindent}{0pt}
\setlength{\parskip}{5pt}
\setlength{\emergencystretch}{3em}
\newcommand{\N}{\mathbb{N}}
\newcommand{\Z}{\mathbb{Z}}
\newcommand{\Q}{\mathbb{Q}}
\newcommand{\R}{\mathbb{R}}
\newcommand{\C}{\mathbb{C}}
\newcommand{\F}{\mathbb{F}}
\newcommand{\A}{\mathbb{A}}
\newcommand{\PP}{\mathbb{P}}
\newcommand{\E}{\mathbb{E}}
\newcommand{\eps}{\varepsilon}
\newcommand{\dd}{\,\mathrm{d}}
\begin{document}
\section*{UnsolvedMath ID: 800325}
\textbf{Source problem code:} catalogue entry
\subsection{Definitions and mathematical statement}
Let \(G\) be a group and \(S\subseteq G\). The set \(S\) is symmetric if \(S=S^{-1}\), where \(S^{-1}=\{s^{-1}:s\in S\}\). The Cayley graph \(\Gamma(G,S)\) has vertex set \(G\), with an edge from \(g\) to \(gs\) for each \(g\in G\) and \(s\in S\). For symmetric \(S\), these adjacencies are reversible. The supplied fragment does not specify whether \(S\) must be finite or generate \(G\), or how loops and edge multiplicities are treated.

For clarity about the terminology suggested by the context, fix a field \(K\) of characteristic zero. A formal power series \(A(z)\in K[[z]]\) is differentially algebraic over \(K(z)\) if there exist an integer \(r\geq 0\) and a nonzero polynomial \(P\in K(z)[Y_0,\ldots,Y_r]\) such that
\[
P\bigl(A(z),A'(z),\ldots,A^{(r)}(z)\bigr)=0.
\]
It is differentially transcendental over \(K(z)\) if no such \(r\) and \(P\) exist. The source uses the terms ``D-algebraic sequence'' and ``D-transcendental Cayley graph'' without defining the associated series or its coefficient field in the supplied record; the preceding definition therefore explains terminology without resolving those missing choices.

\textbf{Problem fragment (Pak, 2018, Problem 9).} Find an explicit construction of \(\Gamma(G,S)\) with \(S\) symmetric. The supplied statement ends immediately after ``and \(\$\)'', leaving an additional required condition missing. Its surrounding discussion suggests that the intended construction should yield a D-transcendental Cayley graph, but this interpretation cannot replace the missing text.

The source introduces the quantity
\[
\rho(G,S):=\lim_{n\to\infty}a_n(G,S)^{1/n},
\]
calling it a spectral radius. It does not define \(a_n(G,S)\) or supply hypotheses ensuring that this limit exists. Consequently, neither a counting interpretation for \(a_n(G,S)\) nor a normalization of \(\rho(G,S)\) is imposed here. If the intended series is \(\sum_{n\geq0}a_n(G,S)z^n\), that identification still requires confirmation.

The supplied background reports an explicit but highly artificial non-symmetric set \(S\subset F_k\times F_\ell\) in a then-forthcoming paper labelled [GP-ADE], where \(F_k\) and \(F_\ell\) denote free groups of the indicated, unspecified ranks. It also reports that [KP2], using tools from [KP1], establishes an uncountable set of spectral radii for groups, and argues from countability of D-algebraic sequences that D-transcendental Cayley graphs with symmetric \(S\) exist nonconstructively. These are supplied contextual claims; their hypotheses, coefficient conventions, and bibliographic identities are unavailable in this record.
\subsection{Short English statement}
Pak's supplied Problem 9 asks for an explicit Cayley graph with a symmetric defining set. The context suggests a D-transcendence requirement, but the statement is truncated before its final condition, and its counting sequence is undefined. The registry labels the problem open; that label and the contextual existence claims are supplied information, not independently verified mathematical assessments.
\subsection{Sources}
\begin{itemize}
\item Local UnsolvedMath registry, permanent ID \(800325\), exact source problem code \texttt{catalogue entry}, title ``Complexity in Enumerative Combinatorics --- Problem 9''. The supplied record belongs to the AMR Open Problem Lists collection and attributes the problem to Pak (2018). Local registry locations are \texttt{data/unsolvedmath/problems.json}, \texttt{data/unsolvedmath/additions/}, and \texttt{catalog/problems.json}; the particular containing file was not independently inspected.
\item Supplied source attribution: Pak, \emph{Complexity Problems in Enumerative Combinatorics} (2018), Problem 9, \url{https://arxiv.org/abs/1803.06636}. The record gives access date 2026-10-16 and extraction method \texttt{arxiv\_latex}. These are supplied provenance details; no source visit or independent verification was performed.
\item The statement cites [GP-ADE], described there as forthcoming, and [KP2] and [KP1]. No bibliographic entries or links identifying these three references are supplied. The word ``forthcoming'' is retained as historical source wording, not a current publication-status assertion.
\item The registry supplies status \texttt{open}, creation date 2026-10-16, update date 2026-10-16, and upstream revision \texttt{d371ee917ac9647fbb06cd6fc529deec6f3dda82}. These metadata do not establish mathematical correctness or independently verified current status.
\end{itemize}
\end{document}
